\section{Related Work}
\label{sec:related}

\para{LLMs as Multiverse Simulators.} The conceptualization of Large Language Models as simulators rather than agents was first formalized by \citet{reynolds2021multiversal}, who introduced the idea of "multiversal views" on language models. They argued that LLMs represent a superposition of possible futures, which they termed "dynamic multiplicity." This was further refined by \citet{shanahan2023roleplay}, who described LLMs as "simulacra" that inhabit characters through role-play. Unlike these theoretical works, we provide an empirical framework to quantify this multiplicity through \method and characterize the entropy of the resulting distributions.

\para{Generative Agents and Social Simulation.} The "Smallville" experiment by \citet{park2023generative} demonstrated that societies of LLM-driven agents can exhibit emergent social behaviors through memory and planning architectures. Recent work like \textsc{AgentSociety} \cite{piao2025agentsociety} has scaled this to city-level simulations. While these studies focus on the *believability* of agent societies, our work focuses on the *breadth* of the simulated outcomes. We are less interested in whether an agent's day-to-day routine is realistic and more interested in whether the model can capture the full distribution of possible human responses to a specific crisis or conflict.

\para{LLMs in Social Science and Economics.} \citet{horton2023homo} proposed the "Homo Silicus" framework, arguing that LLMs can serve as computational models of human behavior for economic experiments. Subsequent work has used LLMs to simulate voting behavior, negotiation, and social norm adherence \cite{argyle2023out, chwang2023simulating}. Our work builds on this by identifying a critical limitation: "alignment-induced mode collapse." We show that standard models like \gpt exhibit a strong central tendency toward "safe" outcomes, which may bias social science findings if not properly accounted for via persona conditioning.

\para{Branching and Diversity in Generation.} Exploring the diversity of LLM outputs has been studied in the context of creative writing and code generation \cite{wang2022self, li2022competition}. Techniques like "Best-of-N" sampling or high-temperature generation are common ways to explore the output space. We extend these techniques to behavioral simulation, using a systematic branching framework to map the manifold of social strategy.
